\bindcontent{
\section*{Study Topic 3: Probability Theory}

\begin{concept}{Basic Probability}
    Marginal probability is the likelihood of a single event occurring.
    \begin{formula}{Classical Probability}
        $\displaystyle P(A) = \frac{n(A)}{N}$
    \end{formula}
\end{concept}

\begin{concept}{Probability Rules}
    The \textbf{General Addition Rule} calculates the probability of event A or event B occurring.
    \begin{formula}{General Addition Rule}
        $P(A \cup B) = P(A) + P(B) - P(A \cap B)$
    \end{formula}
\end{concept}

\begin{concept}{Conditional Probability}
    The probability of event B given that event A has occurred.
    \begin{formula}{Conditional Probability}
        $\displaystyle P(B | A) = \frac{P(A \cap B)}{P(A)}$
    \end{formula}
\end{concept}

\begin{concept}{Independence of Events}
    Two events are independent if $P(B | A) = P(B)$ or:
    \begin{formula}{Multiplication Rule for Independent Events}
        $P(A \cap B) = P(A) \times P(B)$
    \end{formula}
\end{concept}

\begin{concept}{Multiplication Rule}
    Used for a sequence of dependent events (e.g., without replacement).
    \begin{formula}{General Multiplication Rule}
        $P(A \cap B \cap C) = P(A) \times P(B|A) \times P(C|A \cap B)$
    \end{formula}
\end{concept}
}
